\subsubsection{SHAPformer: Explainable Time-Series Forecasting with Sampling-Free SHAP for Transformers}

\paragraph{Ground-truth construction.}
\citet{hertel2026shapformer} construct a synthetic electricity-load forecasting
benchmark. Each example represents the task of predicting the hourly
electricity demand for the following week: the known data-generating function
$g$ returns a 168-hour load curve. Its inputs include a load profile composed
of separate daily patterns for workdays, Saturdays, and Sundays or holidays;
calendar information such as month, hour, weekday, and holiday status; a
temperature curve; and noise components. The calendar variables place the
appropriate daily pattern along the future week, while temperature and the
remaining components modify the resulting demand. The same electricity-load
generator is treated as the function to be explained by SHAP. Therefore, the
reference explanation states how much each input component contributes to the
generated electrical load at each hour of the forecast week. In the authors'
generation arrays, positions $0$--$167$ form the observed week and positions
$168$--$335$ form the future week. With the common origin placed at position
167, the final observed load is at $\tau=0$, the future value stored at
position 168 is $h=1$, and the first element of the 168-dimensional output is
therefore the first unseen hour.

For a coalition of known components $S$, the absent components are repeatedly
resampled while respecting the dependencies encoded in the generation process.
The corresponding horizon-specific coalition value is estimated by the
Monte Carlo average represented as

\begin{equation}
    v_{g,h}(S;x)
    =
    \mathbb{E}_{X_{\bar S}\sim q(\cdot\mid x_S)}
    \!\left[g_h\!\left(x_S,X_{\bar S}\right)\right],
\end{equation}

where $x_S$ contains the component values retained from the example and
$X_{\bar S}$ represents alternative values sampled for the absent components.
Here $q$ names the conditional resampling law implemented by the authors'
dependency-aware procedure; the expectation is approximated with repeated
samples from that procedure.
Shapley values are subsequently calculated from the marginal differences
$v_{g,h}(S\cup\{i\};x)-v_{g,h}(S;x)$. This construction yields signed, local
ground-truth contributions for every feature group and every step of the
forecast horizon. At forecast step $h$, $\phi_{i,h}(x)>0$ means that, averaged
over the coalition contexts in the Shapley sum, inserting the observed value
of component $i$ increases the expected load generated at $h$; conversely,
$\phi_{i,h}(x)<0$ means that it decreases the expected generated load. Thus, the
sign records the direction of the component's contribution relative to the
resampling reference, whereas $\lvert\phi_{i,h}(x)\rvert$ records its strength.
By Shapley efficiency, these signed contributions reconstruct the example
relative to the empty-coalition baseline,

\begin{equation}
    v_{g,h}(\varnothing;x)
    + \sum_{i=1}^{P}\phi_{i,h}(x)
    = g_h(x).
\end{equation}

These reference contributions can then be compared with the explanations
produced for a trained forecasting model.

\paragraph{SHAPformer evaluation.}
The synthetic validation compares SHAPformer's learned explanations with the
generator-based SHAP values at three levels. First, local explanations from
$m$ test samples are converted into a global importance percentage for feature
group $i$ by summing absolute SHAP values over samples and forecast steps,

\begin{equation}
    G_i
    = 100\,
    \frac{\sum_{n=1}^{m}\sum_{h=1}^{H}
    \left|\phi_{i,h}\!\left(x^{(n)}\right)\right|}
    {\sum_{i'=1}^{P}\sum_{n=1}^{m}\sum_{h=1}^{H}
    \left|\phi_{i',h}\!\left(x^{(n)}\right)\right|},
\end{equation}

and the resulting feature-importance bars are placed beside the corresponding
ground-truth bars. Because this aggregation uses absolute values, the bars
summarize total contribution strength after positive and negative effects have
been combined; dependence plots and local explanation curves retain the sign
and therefore preserve the direction of each effect. Second, dependence plots
compare feature values with their
SHAP values for the 24-hour forecast step and use color to indicate the
strongest interacting variable. The SHAPformer plots are examined alongside
the generator plots to determine whether known daily, weekly, holiday, month,
and multiplier effects are recovered. Third, selected local explanation curves
are compared with ground-truth curves in Appendix C. The paper reports close
agreement for most features, successful suppression of the two noise features,
and an underestimated month effect. Consequently, SHAPformer supplies
numerical ground-truth attribution values, while the published explanation
validation judges their proximity primarily through global bars, dependence
plots, and local-curve inspection rather than a single automatic agreement
metric. Forecast RMSE and explanation inference time are reported separately as
quantitative predictive-performance and computational-efficiency measures.

